\scanpage{111}

\uline{Proposition} \quad Soit $f(z)$ une fonction entière \deleted{\unsure{satisfaisant}} \deleted{\unsure{majorée}} \inserted{dont} \deleted{\unsure{avec|dont}} \inserted{les} coefficients sont majorés en module par ceux d'une fonction de genre $0$ $h(z)$. Si de plus les racines non nulles $z_i$ de $f(z)$ sont telles que $\sum\frac{1}{|z_i|}<+\infty$, $f(z)$ est de genre zéro.

\deleted{En effet, $h(z)$ étant d'}

Démonstration. On peut supposer que $h(z)$ est de la forme $k^{\mathrm{te}}\times\sum_{i=1}^{\infty}(1+\alpha_i z)$, où les $\alpha_i$ sont $>0$, et $\sum\alpha_i<+\infty$. Comme $h(z)$ est d'\uline{ordre} $\leq1$, il en est de $=$ de $f(z)$, donc on aura,
\[
f(z)=e^{\alpha+\beta z}z^m\prod_i\left(1-\frac{z}{z_i}\right)
\]

Il faut montrer que $\beta=0$. On peut supposer pour simplifier $\beta>0$, sinon on remplacerait $f(z)$ par $f(\Theta z)$, avec $\Theta=\beta/|\beta|$. On a alors $\displaystyle |f(z)|=\unsure{K_1^{\mathrm{te}}|K^{\mathrm{te}}}\times e^{\beta x}|z|^m\prod\left|1-\frac{z}{z_i}\right|$\unsure{.}

D'autre part, on sait que \deleted{\unsure{si une}} \inserted{toute} fonction de genre zéro, telle $h(z)$, admet une majoration de la forme $|h(z)|\leq\mathrm C^{\mathrm{te}}\times e^{\varepsilon|z|}$, où on peut prendre $\varepsilon$ arbitraire. \deleted{Si alors on prend $\varepsilon\unsure{<|\leq}\beta$, on aura}
\[
\deleted{\deleted{|f(z)|\leq}|z|^m\prod\left|1-\frac{z}{z_i}\right|\leq\unsure{K|C}e^{\varepsilon|z|}}
\]
\deleted{\unsure{avec|pour|tout} $z$}
\margin{112}
